% paper_11_prolate_spheroidal.tex — GeoVac Paper 11
% The Molecular Fock Projection: Diatomic Molecules as
% Prolate Spheroidal Lattices
% Status: Draft, March 12, 2026
% Target: arXiv (physics.comp-ph / quant-ph)
\documentclass[aps,pra,twocolumn,superscriptaddress,longbibliography]{revtex4-2}

\usepackage{amsmath,amssymb,amsthm,graphicx,xcolor,bm,braket,dcolumn,booktabs}
\usepackage{hyperref}

\newtheorem{theorem}{Theorem}
\newtheorem{corollary}{Corollary}
\newtheorem{proposition}{Proposition}

\begin{document}

\title{The Molecular Fock Projection: Diatomic Molecules as\\
Prolate Spheroidal Lattices}
\thanks{Paper 11 in the Geometric Vacuum series}

\author{J.~Loutey}
\affiliation{Independent Researcher, Kent, Washington}

\date{March 12, 2026}

%% ========================================================================
%% ABSTRACT
%% ========================================================================

\begin{abstract}
Fock's 1935 stereographic projection maps the hydrogen atom onto the
unit three-sphere $S^3$, where the Schr\"odinger equation becomes a
pure Laplace--Beltrami eigenvalue problem with no potential term.  The
GeoVac program discretizes this $S^3$ as a graph Laplacian, achieving
$<0.1\%$ error for one-electron atoms with zero free parameters
(Papers~0, 1, 7).  For molecules, however, $S^3$ is not the natural
geometry: the single-$p_0$ momentum scale cannot encode the
$R$-dependent bonding physics (Paper~8--9, negative theorem).

We resolve this by identifying the \emph{natural geometry principle}:
every separable quantum system has a natural lattice---the lattice
that discretizes the Laplace--Beltrami operator in the coordinate
system where separation occurs.  For one-center systems that geometry
is $S^3$; for two-center systems it is the prolate spheroid.  We
discretize the separated prolate spheroidal equations for H$_2^+$
using a spectral angular solver (Legendre basis) and two radial
solvers---a self-adjoint finite-difference (FD) scheme and a spectral
Laguerre basis---with Brent root-finding in the separation
parameter $c^2 = -R^2 E_{\mathrm{elec}}/2$ to enforce self-consistency.
For $\sigma$ states ($m = 0$), the spectral Laguerre matrix elements
are computed algebraically from three-term recurrence relations,
eliminating numerical quadrature entirely.
The method has \emph{zero} free parameters: all inputs are nuclear
charges, internuclear distance, and basis/grid resolution.

Results for H$_2^+$ with the spectral Laguerre radial solver
($n_{\mathrm{basis}} = 20$): minimum energy
$E_{\min} = -0.6025$~Ha (0.0002\% error vs.\ exact $-0.6026$~Ha),
equilibrium bond length $R_{\mathrm{eq}} = 2.005$~bohr
(0.38\% error), and correct dissociation to H~+~p at $R \to \infty$.
The spectral solver achieves a $250\times$ dimension reduction,
$270\times$ speedup, and $5000\times$ accuracy improvement over
the FD solver, whose 1.01\% error was a discretization artifact.
This establishes the prolate spheroidal lattice as the molecular
analog of Fock's atomic $S^3$ projection.
\end{abstract}

\maketitle

%% ========================================================================
%% I. INTRODUCTION
%% ========================================================================

\section{Introduction}
\label{sec:intro}

The Geometric Vacuum program rests on a simple observation: the
hydrogen atom's Schr\"odinger equation, when projected onto momentum
space via Fock's 1935 conformal map~\cite{Fock1935}, becomes the
eigenvalue problem for the Laplace--Beltrami operator on the unit
three-sphere $S^3$.  The Coulomb potential $-Z/r$ is not an input
force law but a coordinate distortion produced by the stereographic
projection from $S^3$ to flat $\mathbb{R}^3$.  This equivalence,
established symbolically with 18 independent proofs~\cite{loutey_paper7},
underlies the entire GeoVac computational framework: discretize $S^3$
as a graph, solve the graph Laplacian eigenvalue problem, and recover
atomic energies to $<0.1\%$ error with no fitted
parameters~\cite{loutey_paper0, loutey_paper1}.

The extension to molecules has proven more subtle.  Paper~8--9
established the \emph{bond sphere} picture---a diatomic molecule
as a single $S^3$ with two weighted poles---and derived the $SO(4)$
Wigner $D$-matrix elements that rotate atomic orbitals from one
center to another~\cite{loutey_paper8}.  However, that same paper
proved a negative structural theorem: in the molecular Sturmian basis
with a shared momentum scale $p_0$, the Hamiltonian is proportional
to the overlap matrix ($H_{ij} \propto S_{ij}$), rendering the
generalized eigenvalues independent of $R$.  Binding in this
framework requires the full $R$-dependent eigenvalue spectrum
$\beta_k(R)$ from the continuous prolate spheroidal PDE---reintroducing
the very continuous spatial geometry the graph was meant to replace.

For heteronuclear systems, the situation is worse: no shared $p_0$
exists that simultaneously represents both atomic length scales.
The LCAO approach (atom-centered lattices with bridges) partially
circumvents this, achieving $D_e^{\text{CP}} = 0.093$~Ha for LiH
at nmax$= 3$ (1.0\% error)~\cite{loutey_lcao_fci}, but carries
intrinsic BSSE and an $R$-independent kinetic energy that shifts
$R_{\mathrm{eq}}$ to $\sim 2.5$~bohr (expt.\ 3.015).

This paper takes a fundamentally different approach.  Rather than
forcing molecules onto $S^3$, we ask: \emph{what is the natural
geometry for two-center systems?}  The answer is immediate from
classical mathematical physics: the two-center Coulomb problem
separates in prolate spheroidal coordinates $(\xi, \eta, \varphi)$,
just as the one-center problem separates in spherical
coordinates~\cite{Bates1953, Morse1953}.  Fock's projection exploits
spherical separation via $SO(4)$ symmetry; the molecular analog
exploits prolate spheroidal separation.

We formalize this as the \textbf{natural geometry principle}:
\begin{quote}
\emph{Every separable quantum system has a natural lattice---the
lattice that discretizes the Laplace--Beltrami operator in the
coordinate system where separation occurs.}
\end{quote}
For atoms, the natural lattice is the $S^3$ graph.  For diatomics,
it is the prolate spheroidal lattice.  The GeoVac program is not
about $S^3$ specifically---it is about identifying and discretizing
the natural conformal geometry for each class of quantum systems.

The remainder of this paper implements this principle for H$_2^+$,
the simplest diatomic.  Section~\ref{sec:coords} reviews prolate
spheroidal coordinates.  Section~\ref{sec:separation} derives the
separated equations.  Sections~\ref{sec:angular} and \ref{sec:radial}
describe the angular (spectral) and radial (finite-difference) solvers.
Section~\ref{sec:spectral_radial} introduces the spectral Laguerre
radial solver that achieves $0.0002\%$ accuracy.
Section~\ref{sec:rootfinding} explains the self-consistency
root-finding.  Section~\ref{sec:results} presents H$_2^+$ results:
PES, spectroscopic constants, and grid convergence.
Section~\ref{sec:discussion} discusses implications for the GeoVac
program.  Section~\ref{sec:conclusion} concludes.

%% ========================================================================
%% II. PROLATE SPHEROIDAL COORDINATES
%% ========================================================================

\section{Prolate Spheroidal Coordinates}
\label{sec:coords}

Place two nuclei $A$ and $B$ at the foci of a prolate spheroid,
separated by distance $R$.  The prolate spheroidal coordinates
$(\xi, \eta, \varphi)$ are defined by
\begin{align}
  \xi  &= \frac{r_A + r_B}{R}, & \xi &\in [1, \infty), \label{eq:xi_def} \\
  \eta &= \frac{r_A - r_B}{R}, & \eta &\in [-1, +1], \label{eq:eta_def} \\
  \varphi &= \arctan(y/x),     & \varphi &\in [0, 2\pi), \label{eq:phi_def}
\end{align}
where $r_A$ and $r_B$ are distances from the electron to nuclei $A$
and $B$.  Surfaces of constant $\xi$ are confocal ellipsoids (with
$\xi = 1$ the degenerate line segment joining the foci); surfaces
of constant $\eta$ are confocal hyperboloids ($\eta = \pm 1$ are
the two half-axes beyond the nuclei).

The metric tensor has diagonal elements
\begin{equation}
  g_{\xi\xi} = \frac{R^2}{4} \frac{\xi^2 - \eta^2}{\xi^2 - 1},
  \quad
  g_{\eta\eta} = \frac{R^2}{4} \frac{\xi^2 - \eta^2}{1 - \eta^2},
  \quad
  g_{\varphi\varphi} = \frac{R^2}{4}(\xi^2 - 1)(1 - \eta^2),
  \label{eq:metric}
\end{equation}
and the volume element is
\begin{equation}
  dV = \frac{R^3}{8}(\xi^2 - \eta^2)\,d\xi\,d\eta\,d\varphi.
  \label{eq:volume}
\end{equation}

The nuclear distances take the simple forms
\begin{equation}
  r_A = \frac{R}{2}(\xi + \eta), \qquad
  r_B = \frac{R}{2}(\xi - \eta),
  \label{eq:nuclear_distances}
\end{equation}
so the nuclear attraction potential for charges $Z_A$ and $Z_B$ is
\begin{equation}
  V = -\frac{Z_A}{r_A} - \frac{Z_B}{r_B}
    = -\frac{2}{R}\left[
        \frac{(Z_A + Z_B)\xi + (Z_B - Z_A)\eta}{\xi^2 - \eta^2}
      \right].
  \label{eq:potential}
\end{equation}

The Laplacian in prolate spheroidal coordinates is
\begin{equation}
  \nabla^2 = \frac{4}{R^2(\xi^2 - \eta^2)}
  \left[
    \frac{\partial}{\partial\xi}\!\left((\xi^2\!-\!1)\frac{\partial}{\partial\xi}\right)
    + \frac{\partial}{\partial\eta}\!\left((1\!-\!\eta^2)\frac{\partial}{\partial\eta}\right)
    + \frac{\xi^2\!-\!\eta^2}{(\xi^2\!-\!1)(1\!-\!\eta^2)}\frac{\partial^2}{\partial\varphi^2}
  \right].
  \label{eq:laplacian}
\end{equation}

%% ========================================================================
%% III. SEPARATION OF THE TWO-CENTER COULOMB PROBLEM
%% ========================================================================

\section{Separation of the Two-Center Coulomb Problem}
\label{sec:separation}

The one-electron Schr\"odinger equation
$(-\tfrac{1}{2}\nabla^2 + V)\Psi = E_{\mathrm{elec}}\Psi$
with the potential~\eqref{eq:potential} can be written, after
multiplying both sides by $-R^2(\xi^2 - \eta^2)/4$, as
\begin{equation}
  \left[\hat{L}_\xi + \hat{L}_\eta
    + \frac{\xi^2 - \eta^2}{(\xi^2-1)(1-\eta^2)}
      \frac{\partial^2}{\partial\varphi^2}\right]\Psi
  = 0,
  \label{eq:full_pde}
\end{equation}
where we define the $\xi$- and $\eta$-operators
\begin{align}
  \hat{L}_\xi &= \frac{\partial}{\partial\xi}
    \!\left[(\xi^2 - 1)\frac{\partial}{\partial\xi}\right]
    + a\xi - c^2\xi^2, \label{eq:L_xi} \\
  \hat{L}_\eta &= \frac{\partial}{\partial\eta}
    \!\left[(1 - \eta^2)\frac{\partial}{\partial\eta}\right]
    + b\eta + c^2\eta^2, \label{eq:L_eta}
\end{align}
with the dimensionless parameters
\begin{equation}
  a = R(Z_A + Z_B), \quad
  b = R(Z_B - Z_A), \quad
  c^2 = -\frac{R^2 E_{\mathrm{elec}}}{2}.
  \label{eq:params}
\end{equation}
For bound states ($E_{\mathrm{elec}} < 0$), $c^2 > 0$.

The ansatz $\Psi(\xi,\eta,\varphi) = F(\xi)\,G(\eta)\,e^{im\varphi}$
separates Eq.~\eqref{eq:full_pde} into two ordinary differential
equations coupled through a separation constant $A$:

\paragraph{$\xi$-equation (radial):}
\begin{equation}
  \frac{d}{d\xi}\!\left[(\xi^2 - 1)\frac{dF}{d\xi}\right]
  + \left(A + a\xi - c^2\xi^2 - \frac{m^2}{\xi^2 - 1}\right)F = 0.
  \label{eq:xi_equation}
\end{equation}

\paragraph{$\eta$-equation (angular):}
\begin{equation}
  \frac{d}{d\eta}\!\left[(1 - \eta^2)\frac{dG}{d\eta}\right]
  + \left(-A + c^2\eta^2 + b\eta - \frac{m^2}{1 - \eta^2}\right)G = 0.
  \label{eq:eta_equation}
\end{equation}

The angular equation~\eqref{eq:eta_equation} is a generalization of
the associated Legendre equation with additional $c^2\eta^2$ and
$b\eta$ terms.  For homonuclear diatomics ($Z_A = Z_B$), $b = 0$ and
the $\eta$-equation reduces to the spheroidal wave equation.  The
radial equation~\eqref{eq:xi_equation} has an irregular singular
point at $\xi = 1$ (the line segment between the nuclei) and an
irregular singular point at $\xi \to \infty$.

The key feature of the separated system is that both
equations~\eqref{eq:xi_equation} and \eqref{eq:eta_equation}
contain the same three unknowns: the energy $E_{\mathrm{elec}}$
(via $c^2$), the separation constant $A$, and the quantum numbers
$(m, n_\xi, n_\eta)$.  For fixed $m$ and state indices, the problem
reduces to finding the value of $c^2$ for which both equations yield
the same $A$.

%% ========================================================================
%% IV. THE ANGULAR EQUATION: SPECTRAL SOLVER
%% ========================================================================

\section{The Angular Equation: Spectral Solver}
\label{sec:angular}

The angular equation~\eqref{eq:eta_equation} is an eigenvalue problem
for $A$ at fixed $c^2$.  We solve it using a spectral method with
Legendre polynomial basis functions~\cite{Flammer1957}.

\subsection{Legendre expansion}

Expand $G(\eta) = \sum_{l=|m|}^{l_{\max}} a_l P_l^{|m|}(\eta)$
in associated Legendre polynomials.  Substituting into
Eq.~\eqref{eq:eta_equation} and using the recurrence relations
for $P_l^m$ produces a symmetric tridiagonal (pentadiagonal for
$b \neq 0$) generalized eigenvalue problem
\begin{equation}
  \mathbf{M}\,\mathbf{a} = A\,\mathbf{a},
  \label{eq:angular_eigenproblem}
\end{equation}
where the matrix $\mathbf{M}$ has elements determined by the
three-term recurrence for $\eta^2 P_l^m$ and (when $b \neq 0$)
$\eta P_l^m$.

For homonuclear diatomics ($b = 0$), the matrix $\mathbf{M}$ is
symmetric and tridiagonal, and the eigenvalues converge exponentially
in $l_{\max}$.  With $l_{\max} = 50$, the ground-state separation
constant $A$ is converged to machine precision for all $c^2$ values
of interest.

\subsection{Boundary conditions}

The associated Legendre polynomials automatically satisfy the correct
boundary conditions: regularity at $\eta = \pm 1$.  No additional
boundary conditions need to be imposed, because the factor
$(1 - \eta^2) \to 0$ at the endpoints provides a natural zero-flux
condition for the self-adjoint operator.

For homonuclear diatomics, the solutions have definite parity under
$\eta \to -\eta$: even parity ($\sigma_g$, $\pi_u$, \ldots) or odd
parity ($\sigma_u$, $\pi_g$, \ldots).  This parity is automatically
captured by the Legendre expansion through the selection of even or
odd $l$ values.

%% ========================================================================
%% V. THE RADIAL EQUATION: SELF-ADJOINT FINITE DIFFERENCES
%% ========================================================================

\section{The Radial Equation: Self-Adjoint Finite Differences}
\label{sec:radial}

The radial equation~\eqref{eq:xi_equation} is solved numerically
on a uniform grid in $\xi$.  The critical design choice is the
discretization of the self-adjoint differential operator
$d/d\xi\,[p(\xi)\,dF/d\xi]$ where $p(\xi) = \xi^2 - 1$.

\subsection{Self-adjoint stencil}

The standard second-order self-adjoint stencil~\cite{Trefethen2000}
for $d/d\xi\,[p(\xi)\,dF/d\xi]$ on a uniform grid with spacing $h$
evaluates $p$ at the half-points:
\begin{equation}
  \frac{1}{h^2}\left[
    p_{i+1/2}(F_{i+1} - F_i) - p_{i-1/2}(F_i - F_{i-1})
  \right],
  \label{eq:sa_stencil}
\end{equation}
where $p_{i\pm 1/2} = (\xi_i \pm h/2)^2 - 1$.  This yields a
symmetric tridiagonal matrix with diagonal elements
\begin{equation}
  d_i = -\frac{p_{i+1/2} + p_{i-1/2}}{h^2}
        + A + a\xi_i - c^2\xi_i^2 - \frac{m^2}{\xi_i^2 - 1}
  \label{eq:diag}
\end{equation}
and off-diagonal elements
\begin{equation}
  e_i = \frac{p_{i+1/2}}{h^2}.
  \label{eq:offdiag}
\end{equation}

\subsection{Boundary conditions}

Two boundary conditions are required:

\paragraph{At $\xi = 1$ (nuclear axis):}
For $\sigma$ states ($m = 0$), the wavefunction is nonzero on the
internuclear axis.  Since $p(1) = 0$, the self-adjoint operator has
a natural \emph{Neumann} boundary condition at $\xi = 1$: the flux
$p(\xi)\,dF/d\xi \to 0$ as $\xi \to 1^+$ regardless of the value
of $dF/d\xi$.  Computationally, this is implemented by setting
$p_{1/2} = 0$ in the stencil~\eqref{eq:sa_stencil}, which modifies
the first diagonal element to $d_1 = -p_{3/2}/h^2 + q_1$.

For $\pi$ states ($m \neq 0$), the wavefunction vanishes on the axis
by the $e^{im\varphi}$ factor, and a Dirichlet condition $F(1) = 0$
is appropriate.

\paragraph{At $\xi = \xi_{\max}$ (asymptotic region):}
A Dirichlet condition $F(\xi_{\max}) = 0$ is imposed, with
$\xi_{\max}$ chosen large enough that the bound-state wavefunction
has decayed to negligible amplitude.  For H$_2^+$ near equilibrium,
$\xi_{\max} = 25$ is sufficient; for large $R$, $\xi_{\max}$ is
scaled with $R$.

\subsection{Eigenvalue interpretation}

With $A$ fixed from the angular solver, the radial
operator~\eqref{eq:xi_equation} is a Sturm--Liouville problem
parameterized by $c^2$.  At the \emph{correct} $c^2$, the largest
eigenvalue of the discretized operator equals zero---the bound state
sits exactly at the top of the spectrum.  At other values of $c^2$,
the top eigenvalue is nonzero.  This provides a scalar residual
function $f(c^2) = \lambda_{\max}(c^2)$ for root-finding
(Section~\ref{sec:rootfinding}).

The grid offset $\xi_{\min} = 1 + \delta$ (with $\delta = 5 \times
10^{-4}$) avoids the singular point $\xi = 1$ while keeping the
first grid point close enough to capture the near-nuclear wavefunction.

%% ========================================================================
%% V.B SPECTRAL LAGUERRE RADIAL SOLVER
%% ========================================================================

\section{Spectral Laguerre Radial Solver}
\label{sec:spectral_radial}

The finite-difference radial solver of Sec.~\ref{sec:radial} converges
as $O(1/N_\xi)$, requiring $N_\xi = 5000$ grid points for sub-percent
accuracy.  This section describes a spectral alternative that achieves
$5000\times$ better accuracy with $250\times$ fewer degrees of freedom.

\subsection{Laguerre basis}

The radial equation~\eqref{eq:xi_equation} is defined on
$\xi \in [1, \infty)$ with a regular singular point at $\xi = 1$
(where $p(\xi) = \xi^2 - 1 \to 0$) and exponential decay as
$\xi \to \infty$.  These features suggest a Laguerre basis:
\begin{equation}
  \varphi_n(\xi) = e^{-\alpha(\xi - 1)}\,L_n\bigl(2\alpha(\xi - 1)\bigr),
  \quad n = 0, 1, \ldots, N_b - 1,
  \label{eq:laguerre_basis}
\end{equation}
where $L_n$ is the Laguerre polynomial and the decay parameter
$\alpha = \sqrt{c^2}$ is adapted to the wavefunction's asymptotic
decay rate at each value of $c^2$.

\subsection{Matrix elements}

The kinetic operator $d/d\xi\,[(\xi^2 - 1)\,dF/d\xi]$ is treated
via integration by parts:
\begin{equation}
  \langle \varphi_m | \hat{T} | \varphi_n \rangle
  = -\int_1^\infty (\xi^2 - 1)\,\varphi_m'(\xi)\,\varphi_n'(\xi)\,d\xi,
  \label{eq:ibp}
\end{equation}
where the boundary terms vanish: $(\xi^2 - 1) = 0$ at $\xi = 1$
and the exponential decay ensures convergence at $\xi \to \infty$.
The potential matrix elements
$\langle \varphi_m | q(\xi) | \varphi_n \rangle$ with
$q(\xi) = A + a\xi - c^2\xi^2 - m^2/(\xi^2 - 1)$ and the overlap
matrix $S_{mn} = \langle \varphi_m | \varphi_n \rangle$ are computed
via Gauss--Laguerre quadrature, which is exact for the polynomial
content of the integrands.  For $m = 0$ ($\sigma$ states), these
matrix elements can alternatively be computed algebraically from
Laguerre moment recurrences, eliminating quadrature entirely
(Sec.~\ref{sec:algebraic}).

\subsection{Generalized eigenvalue problem}

The radial equation becomes the generalized eigenvalue problem
\begin{equation}
  \mathbf{H}\,\mathbf{c} = \lambda\,\mathbf{S}\,\mathbf{c},
  \label{eq:gep}
\end{equation}
where $\mathbf{H} = \mathbf{T} + \mathbf{Q}$ (kinetic plus potential)
and $\mathbf{S}$ is the overlap matrix, both of dimension
$N_b \times N_b$.  The same Brent root-finding loop as in
Sec.~\ref{sec:rootfinding} is used: at each trial $c^2$, the angular
solver provides $A(c^2)$, the spectral radial solver provides
$\lambda_{\max}(c^2, A)$, and the root $\lambda_{\max} = 0$
determines the energy.  The spectral solver is activated via
\texttt{radial\_method='spectral'} with \texttt{n\_basis=20}.

\subsection{Spectral convergence}

Table~\ref{tab:spectral_convergence} compares the finite-difference
and spectral radial solvers for H$_2^+$ at $R = 2.0$~bohr.

\begin{table}[h]
\caption{Comparison of finite-difference and spectral Laguerre
  radial solvers for H$_2^+$ at $R = 2.0$~bohr.
  Exact total energy: $-0.6026$~Ha.
  \label{tab:spectral_convergence}}
\begin{ruledtabular}
\begin{tabular}{ldd}
  \multicolumn{1}{c}{Metric}
  & \multicolumn{1}{c}{FD ($N_\xi = 5000$)}
  & \multicolumn{1}{c}{Spectral ($N_b = 20$)} \\
\hline
  $E_{\mathrm{total}}$ error & 1.01\% & 0.0002\% \\
  $R_{\mathrm{eq}}$ (bohr)   & \sim 2.01 & 2.005 \\
  $R_{\mathrm{eq}}$ error    & 0.65\% & 0.38\% \\
  Wall time (s)               & 4.3 & 0.016 \\
  Matrix dimension            & 5000 & 20 \\
\end{tabular}
\end{ruledtabular}
\end{table}

The spectral solver achieves a $250\times$ dimension reduction,
$270\times$ speedup, and $5000\times$ accuracy improvement over
the FD solver.  The FD solver's 1.01\% error was a discretization
artifact---the $O(1/N_\xi)$ convergence from the singular point
at $\xi = 1$ and the first-order Neumann boundary condition
(Appendix~\ref{app:convergence}) masked the intrinsic accuracy of
the prolate spheroidal basis.  The spectral solver, which treats
the $\xi = 1$ boundary analytically via the vanishing of
$(\xi^2 - 1)$ in the integration-by-parts identity~\eqref{eq:ibp},
reveals that the prolate spheroidal separation itself is far more
accurate than previously reported.

Spectral convergence is rapid: $N_b = 5$ already achieves
sub-$0.01\%$ accuracy, and $N_b = 20$ reaches the
$10^{-4}\%$ level.  The existing FD solver is preserved as the
default for backward compatibility.

\subsection{Algebraic matrix elements for \texorpdfstring{$\sigma$}{sigma} states}
\label{sec:algebraic}

For $m = 0$ ($\sigma$ states), all matrix elements in the spectral
Laguerre solver can be computed \emph{algebraically} from three-term
recurrence relations, eliminating Gauss--Laguerre quadrature entirely.

Define the \emph{Laguerre moment matrices}
\begin{equation}
  M_k[i,j] = \int_0^\infty x^k\,L_i(x)\,L_j(x)\,e^{-x}\,dx,
  \label{eq:laguerre_moment}
\end{equation}
where $L_n$ denotes the Laguerre polynomial of degree~$n$.
The three-term recurrence
\begin{equation}
  x\,L_n(x) = -(n+1)\,L_{n+1}(x) + (2n+1)\,L_n(x) - n\,L_{n-1}(x)
  \label{eq:laguerre_recurrence}
\end{equation}
gives $M_k$ as finite band matrices built recursively from
$M_0 = \mathbf{I}$ (orthonormality): $M_1$ is tridiagonal with
$M_1[n,n] = 2n+1$, $M_1[n, n\pm 1] = -(n+1)$ or $-n$;
$M_2$ is pentadiagonal, and so on.  Higher moments are obtained by
applying the recurrence~\eqref{eq:laguerre_recurrence} to reduce
$M_k$ to a sum of products involving $M_{k-1}$.

Under the substitution $x = 2\alpha(\xi - 1)$, each matrix element
reduces to a finite linear combination of Laguerre moments:
\begin{itemize}
  \item \textbf{Overlap:} $S_{mn} = \delta_{mn}/(2\alpha)$.
  \item \textbf{Potential:} The effective potential
    $q(\xi) = A + a\xi - c^2\xi^2$ (at $m = 0$, no singular term)
    becomes a quadratic in~$x$, giving
    $\mathbf{V} = (2\alpha)^{-1}[q_0\,\mathbf{M}_0 + q_1\,\mathbf{M}_1
    + q_2\,\mathbf{M}_2]$
    where $q_0$, $q_1$, $q_2$ are determined by the coefficients $A$,
    $a$, $c^2$, and $\alpha$.
  \item \textbf{Kinetic:} Writing $\varphi_n' = -2\alpha\,\sum_{j \le n} L_j$
    (from the Laguerre derivative identity) and using the
    integration-by-parts form~\eqref{eq:ibp}, the kinetic matrix becomes
    $\mathbf{K} = -\tfrac{1}{2}\,\mathbf{F}\,\mathbf{M}_1\,\mathbf{F}^T
    - \tfrac{1}{8\alpha}\,\mathbf{F}\,\mathbf{M}_2\,\mathbf{F}^T$,
    where $F_{nj} = 2$ for $j < n$, $F_{nn} = 1$.
\end{itemize}

The resulting overlap and Hamiltonian matrices agree with
Gauss--Laguerre quadrature to machine precision: maximum discrepancy
$4.4 \times 10^{-15}$ for the overlap and $4.1 \times 10^{-12}$ for the
Hamiltonian (the latter reflecting accumulated floating-point arithmetic
in the moment products).  PES energies agree to $< 2 \times 10^{-14}$~Ha
across all bond lengths, confirming exact equivalence.  The algebraic
construction is approximately $1.6\times$ faster than quadrature,
as it avoids computation of Gauss--Laguerre roots and weights.

For $m \neq 0$, the centrifugal term $m^2/(\xi^2 - 1)$ produces
a $1/x$ singularity in the Laguerre variable, and the required
integral $\int_0^\infty L_i\,L_j\,e^{-x}/x\,dx$ diverges at
$x = 0$---it has no finite Laguerre moment expansion.  The solver
therefore falls back to Gauss--Laguerre quadrature for $m \neq 0$
states automatically.  Since the ground state $1s\sigma_g$ and all
$\sigma$ excited states have $m = 0$, the algebraic method covers
the primary use cases of the solver.

%% ========================================================================
%% VI. ROOT-FINDING AND SELF-CONSISTENCY
%% ========================================================================

\section{Root-Finding and Self-Consistency}
\label{sec:rootfinding}

The complete algorithm proceeds as follows:

\begin{enumerate}
  \item \textbf{Input:} $Z_A$, $Z_B$, $R$, grid parameters
        ($N_\xi$, $\xi_{\max}$), azimuthal quantum number $m$.
  \item \textbf{Define residual} $f(c^2)$:
    \begin{enumerate}
      \item Compute $A(c^2)$ from the angular spectral
            solver~(Section~\ref{sec:angular}).
      \item Compute $\lambda_{\max}(c^2, A)$ from the radial FD
            solver~(Section~\ref{sec:radial}).
      \item Return $f = \lambda_{\max}$.
    \end{enumerate}
  \item \textbf{Bracket:} Find $c^2_{\min}$ and $c^2_{\max}$ such
        that $f(c^2_{\min}) > 0$ and $f(c^2_{\max}) < 0$.  For
        homonuclear H$_2^+$, typical brackets are $c^2 \in
        [0.01, R^2(Z_A + Z_B)^2]$.
  \item \textbf{Root-find:} Apply Brent's method~\cite{Brent1973}
        to solve $f(c^2) = 0$ to tolerance $10^{-12}$.
  \item \textbf{Extract energy:}
        $E_{\mathrm{elec}} = -2c^2/R^2$, \quad
        $E_{\mathrm{total}} = E_{\mathrm{elec}} + Z_A Z_B / R$.
\end{enumerate}

The algorithm is completely parameter-free: the only inputs are the
physical quantities ($Z_A$, $Z_B$, $R$) and the numerical resolution
($N_\xi$, $\xi_{\max}$).  No orbital exponents, basis set parameters,
kinetic scaling factors, or bridge coupling constants appear anywhere.

Each evaluation of $f(c^2)$ requires one angular eigenvalue solve
($O(l_{\max})$ for tridiagonal) and one radial eigenvalue solve
($O(N_\xi)$ for tridiagonal).  Brent's method typically converges in
$\sim 50$ iterations, so the total cost per energy point is
$O(50 \times N_\xi) \sim O(N_\xi)$.  For $N_\xi = 5000$, a single
point takes approximately 0.2~s on commodity hardware.

%% ========================================================================
%% VII. RESULTS
%% ========================================================================

\section{Results: \texorpdfstring{H$_2^+$}{H2+} Ground State}
\label{sec:results}

\subsection{Grid convergence}

Table~\ref{tab:convergence} shows the convergence of the ground-state
total energy for H$_2^+$ at $R = 2.0$~bohr as a function of radial
grid size $N_\xi$ (with $\xi_{\max} = 25$).  The exact
value~\cite{Bates1953} is $E_{\mathrm{total}} = -0.6026$~Ha.

\begin{table}[h]
\caption{Grid convergence for H$_2^+$ at $R = 2.0$~bohr.
  The exact total energy is $-0.6026$~Ha.
  \label{tab:convergence}}
\begin{ruledtabular}
\begin{tabular}{rddd}
  \multicolumn{1}{c}{$N_\xi$}
  & \multicolumn{1}{c}{$c^2$}
  & \multicolumn{1}{c}{$E_{\mathrm{total}}$ (Ha)}
  & \multicolumn{1}{c}{Error (\%)} \\
\hline
    500  & 2.114 & -0.5568 & 7.60 \\
   1000  & 2.156 & -0.5779 & 4.11 \\
   2000  & 2.179 & -0.5893 & 2.21 \\
   5000  & 2.193 & -0.5965 & 1.01 \\
  10000  & 2.198 & -0.5990 & 0.59 \\
\end{tabular}
\end{ruledtabular}
\end{table}

The convergence is approximately $O(1/N_\xi)$, consistent with
second-order finite-difference discretization of the radial equation.
Sub-percent accuracy is achieved at $N_\xi = 5000$ (0.2~s wall
time).  The spectral Laguerre radial solver
(Sec.~\ref{sec:spectral_radial}) eliminates this discretization
bottleneck, achieving $0.0002\%$ accuracy with $N_b = 20$.

\subsection{Potential energy surface}

Figure~\ref{fig:pes} shows the H$_2^+$ potential energy surface
computed with $N_\xi = 8000$ over $R \in [1.0, 6.0]$~bohr.  The
curve exhibits a clear minimum near $R = 2.0$~bohr and dissociates
correctly toward $E = -0.5$~Ha (the H atom ground state plus a bare
proton) at large $R$.

\begin{table}[h]
\caption{H$_2^+$ potential energy surface near the minimum
  ($N_\xi = 8000$, $\xi_{\max} = 25$).
  \label{tab:pes}}
\begin{ruledtabular}
\begin{tabular}{dd}
  \multicolumn{1}{c}{$R$ (bohr)}
  & \multicolumn{1}{c}{$E_{\mathrm{total}}$ (Ha)} \\
\hline
  1.50 & -0.5783 \\
  1.60 & -0.5869 \\
  1.70 & -0.5926 \\
  1.80 & -0.5961 \\
  1.90 & -0.5979 \\
  2.00 & -0.5984 \\
  2.10 & -0.5979 \\
  2.20 & -0.5965 \\
  2.50 & -0.5894 \\
  3.00 & -0.5730 \\
  4.00 & -0.5412 \\
  5.00 & -0.5190 \\
  6.00 & -0.5058 \\
\end{tabular}
\end{ruledtabular}
\end{table}

\subsection{Spectroscopic constants}

Table~\ref{tab:spectroscopic} compares the computed spectroscopic
constants with exact values.  Equilibrium geometry and binding
energy are obtained from a quadratic fit to the PES near the minimum.

\begin{table}[h]
\caption{Spectroscopic constants for H$_2^+$ ($N_\xi = 8000$).
  Exact values from Bates \emph{et al.}~\cite{Bates1953}.
  \label{tab:spectroscopic}}
\begin{ruledtabular}
\begin{tabular}{lddd}
  \multicolumn{1}{c}{Quantity}
  & \multicolumn{1}{c}{Computed}
  & \multicolumn{1}{c}{Exact}
  & \multicolumn{1}{c}{Error (\%)} \\
\hline
  $R_{\mathrm{eq}}$ (bohr) & 2.001 & 1.997 & 0.21 \\
  $E_{\min}$ (Ha)          & -0.5984 & -0.6026 & 0.70 \\
  $D_e$ (Ha)               & 0.0984 & 0.1026 & 4.09 \\
  $k$ (Ha/bohr$^2$)        & 0.1032 & \sim 0.10 & \sim 3 \\
\end{tabular}
\end{ruledtabular}
\end{table}

The equilibrium bond length is reproduced to 0.21\% and the minimum
energy to 0.70\%.  The dissociation energy $D_e = -0.5 - E_{\min}$
has a larger relative error (4.09\%) because it is the small
difference between two large numbers.  The force constant $k$
agrees to $\sim 3\%$.

\subsection{Dissociation limit}

At large $R$, the electronic energy approaches $E_{\mathrm{elec}}
\to -1/2$ (the hydrogen atom ground state) and the total energy
approaches $E_{\mathrm{total}} \to -1/2$ (since $V_{NN} = 1/R
\to 0$).  At $R = 10$~bohr with increased grid parameters
($N_\xi = 8000$, $\xi_{\max} = 80$), we obtain $E_{\mathrm{total}}
= -0.476$~Ha, approaching the exact limit of $-0.500$~Ha.  The
slow convergence at large $R$ reflects the need for a large
$\xi_{\max}$ to capture the diffuse wavefunction.

\subsection{Heteronuclear extension: \texorpdfstring{HeH$^{2+}$}{HeH2+}}

The method extends naturally to heteronuclear one-electron systems.
For HeH$^{2+}$ ($Z_A = 2$, $Z_B = 1$) at $R = 2.0$~bohr, we
obtain $E_{\mathrm{total}} = -1.475$~Ha.  The heteronuclear case
introduces a nonzero $b = R(Z_B - Z_A)$ parameter in the angular
equation, breaking the $\eta \to -\eta$ parity symmetry, but the
algorithm handles this without modification.

%% ========================================================================
%% VIII. DISCUSSION
%% ========================================================================

\section{Discussion}
\label{sec:discussion}

\subsection{The natural geometry principle}

The results of this paper, combined with the atomic results of
Papers~0, 1, and 7, establish a pattern:

\begin{itemize}
  \item \textbf{One-center systems (atoms):} The natural geometry
    is the three-sphere $S^3$, accessed via Fock's stereographic
    projection from momentum space.  The Schr\"odinger equation
    becomes the Laplace--Beltrami eigenvalue problem on $S^3$,
    with no potential term.
  \item \textbf{Two-center systems (diatomics):} The natural
    geometry is the prolate spheroid, where the two-center Coulomb
    problem separates into two coupled ODEs.  The bond length $R$
    enters through the metric and the separation parameter $c^2$.
\end{itemize}

In both cases, the strategy is identical: identify the coordinate
system where the Schr\"odinger equation separates, discretize the
resulting Laplace--Beltrami operator, and solve the resulting
eigenvalue problem.  The physics---nuclear charges, distances,
masses---enters only through the coordinate system and its metric.
There are no basis set parameters, no fitted constants, and no
orbital exponents to optimize.

\subsection{Resolution of the bond sphere impasse}

Paper~8--9 proved that the molecular Sturmian basis on a single
$S^3$ with shared momentum scale $p_0$ cannot produce
$R$-dependent binding: the Hamiltonian is proportional to the
overlap matrix, making eigenvalues geometry-independent.  The
prolate spheroidal lattice resolves this impasse by abandoning
the single-center $S^3$ geometry entirely.  Rather than forcing
two centers onto one sphere, each center retains its identity
within the natural two-center geometry.

The bond length $R$ does not enter as an external parameter grafted
onto an $R$-independent topology.  Instead, $R$ is intrinsic to the
prolate spheroidal metric~\eqref{eq:metric}: it sets the interfocal
distance, which determines the coordinate system itself.  This is
analogous to how $p_0$ is intrinsic to the Fock projection for
atoms---it is not a free parameter but is determined self-consistently
by the energy eigenvalue.

\subsection{Comparison with LCAO approach}

The LCAO approach~\cite{loutey_lcao_fci} constructs a molecular
basis by combining atom-centered lattices with bridge elements.
While successful for energetics ($D_e^{\text{CP}} = 0.093$~Ha
for LiH, 1.0\% error), it has inherent limitations:
\begin{enumerate}
  \item \textbf{BSSE:} Combining two lattices creates basis set
    superposition error ($-0.115$~Ha for LiH at nmax$= 3$),
    requiring counterpoise correction.
  \item \textbf{$R$-independent kinetic energy:} The graph Laplacian
    eigenvalues on each atomic lattice do not depend on $R$,
    shifting $R_{\mathrm{eq}}$ to $\sim 2.5$~bohr (expt.\ 3.015).
  \item \textbf{Bridge parameters:} The inter-atomic coupling
    requires STO overlap integrals and Mulliken-type approximations.
\end{enumerate}

The prolate spheroidal lattice has none of these limitations: no
BSSE (single coordinate system), $R$-dependent metric (correct
$R_{\mathrm{eq}}$), and no fitted parameters.  For one-electron
systems, it is clearly superior.

\subsection{Limitations and extensions}

The most important limitations and implemented extensions are:

\paragraph{Two-electron systems (H$_2$).}
We have implemented a prolate spheroidal CI using
single-particle orbitals ($\sigma_g$, $\sigma_u$) from the separated
H$_2^+$ solver and a 2$\times$2 CI in the $^1\Sigma_g^+$ sector.
The $V_{ee}$ matrix elements are computed via azimuthal averaging
of $1/r_{12}$ using the complete elliptic integral $K(k)$ on a
two-dimensional $(\xi, \eta)$ quadrature grid.

With frozen H$_2^+$ orbitals ($Z_{\mathrm{eff}} = 1$), the CI yields
$D_e = 0.077$~Ha (44\% of exact 0.175~Ha) at $R_{\mathrm{eq}} = 1.77$~bohr.
Optimizing the orbital exponent via the Eckart variational method
($Z_{\mathrm{eff}}^* \approx 0.765$ at $R = 1.4$~bohr) reduces
the Coulomb self-repulsion $J_{gg}$ from 0.80 to 0.68~Ha and improves the
result to $D_e = 0.101$~Ha (58\% of exact) with $R_{\mathrm{eq}}$
matching the exact value of 1.401~bohr within fitting precision.

An important structural result is the \emph{pi orbital selection rule}:
$\langle \sigma_g \sigma_g | \pi_+ \pi_- \rangle = 0$ by azimuthal
symmetry, so $\pi$ orbitals cannot improve the $^1\Sigma_g^+$ ground
state via CI mixing.  The 4$\sigma$ CI (6$\times$6 matrix) confirms
this, adding only 0.1~mHa over the 2$\times$2 result.

The remaining 42\% $D_e$ error arises from the minimal two-orbital
CI basis and missing dynamic correlation.  The $V_{ee}$ integrals
converge as $O(1/N_{\mathrm{grid}}^2)$ in grid size.

\paragraph{Heteronuclear systems (HeH$^+$).}
For heteronuclear diatomics, we have implemented independent
per-atom orbital exponent optimization ($Z_{\mathrm{eff},A}$,
$Z_{\mathrm{eff},B}$), which recovers binding in HeH$^+$---a
system that is unbound with frozen HeH$^{2+}$ orbitals due to
excessive Coulomb repulsion ($J_{11} \sim 1.3$~Ha).

\paragraph{Excited states.}
The current solver targets the ground $\sigma_g$ state.  Excited
states ($\sigma_u$, $\pi_g$, $\pi_u$) are accessible by changing
the azimuthal quantum number $m$ and selecting different eigenvalues
of the angular and radial operators.

\paragraph{Higher-order discretization.}
The spectral Laguerre radial solver (Sec.~\ref{sec:spectral_radial})
has replaced the $O(1/N_\xi)$ FD scheme as the high-accuracy option,
achieving $0.0002\%$ error with $N_b = 20$ basis functions---a
$5000\times$ accuracy improvement over the FD solver at $250\times$
smaller matrix dimension.  For $\sigma$ states, all matrix elements
are now algebraic (Sec.~\ref{sec:algebraic}), with no remaining
numerical quadrature in the radial solver.

%% ========================================================================
%% IX. CONCLUSION
%% ========================================================================

\section{Conclusion}
\label{sec:conclusion}

We have demonstrated that the natural geometry for diatomic molecules
is the prolate spheroid, just as the natural geometry for atoms is
the three-sphere $S^3$.  By discretizing the separated prolate
spheroidal equations with a spectral angular solver and a spectral
Laguerre radial solver, we obtain the H$_2^+$ ground-state
energy to $0.0002\%$ accuracy and the equilibrium bond length to
$0.38\%$ accuracy with zero free parameters.  The spectral solver
achieves a $250\times$ dimension reduction and $270\times$ speedup
over the finite-difference radial solver, whose $1.01\%$ error was
a discretization artifact.  For $\sigma$ states, the spectral matrix
elements are computed algebraically from Laguerre moment recurrences,
eliminating numerical quadrature from the radial solver entirely.

This result establishes the \emph{natural geometry principle} as
the organizing idea of the GeoVac program:
\begin{quote}
\emph{Every separable quantum system has a natural lattice---the
lattice that discretizes the Laplace--Beltrami operator in the
coordinate system where separation occurs.}
\end{quote}
The atomic $S^3$ lattice (Papers~0, 1, 7) and the molecular prolate
spheroidal lattice (this paper) are the first two instances of this
principle.  The two-electron extension to H$_2$ via prolate
spheroidal CI achieves 58\% of the exact binding energy with relaxed
orbitals and zero free parameters.  Future work will pursue larger
CI expansions, polyatomic systems with other natural coordinate
systems, and integration of the spectral radial solver into the
full PES scan and composed-geometry pipelines.

The prolate spheroidal lattice also resolves the impasse identified
in Paper~8--9: the single-$p_0$ Sturmian basis on $S^3$ cannot
produce $R$-dependent binding, but the prolate spheroidal coordinate
system encodes $R$ directly in its metric, producing the correct
equilibrium geometry without any fitted parameters.

%% ========================================================================
%% ACKNOWLEDGMENTS
%% ========================================================================

\begin{acknowledgments}
Computational resources were provided by personal hardware.
\end{acknowledgments}

%% ========================================================================
%% APPENDIX A: GRID CONVERGENCE ANALYSIS
%% ========================================================================

\appendix

\section{Grid Convergence Analysis}
\label{app:convergence}

The second-order self-adjoint finite-difference stencil gives
truncation error $O(h^2) = O((\xi_{\max}/N_\xi)^2)$.  Since the
top eigenvalue of the discretized radial operator must equal zero,
the energy error is controlled by the eigenvalue error, which
scales as $O(h^2)$ for the largest eigenvalue.

In practice, the observed convergence rate is closer to $O(1/N_\xi)$
(Table~\ref{tab:convergence}) rather than $O(1/N_\xi^2)$.  This
sub-optimal convergence is attributed to:
\begin{enumerate}
  \item The grid offset $\delta = 5 \times 10^{-4}$ from the
    singular point $\xi = 1$, which introduces an $O(\delta)$
    boundary error.
  \item The Neumann boundary condition implementation, which is
    first-order at the boundary.
  \item The nonuniform behavior of $p(\xi) = \xi^2 - 1$ near
    $\xi = 1$, where $p \to 0$.
\end{enumerate}

The spectral Laguerre solver of Sec.~\ref{sec:spectral_radial}
eliminates all three issues: the integration-by-parts
identity~\eqref{eq:ibp} handles the $\xi = 1$ singularity
analytically, and the exponential basis functions match the
wavefunction's asymptotic decay, achieving $0.0002\%$ accuracy
with $N_b = 20$.

%% ========================================================================
%% APPENDIX B: COMPARISON WITH FAILED APPROACHES
%% ========================================================================

\section{Comparison with Failed Approaches}
\label{app:failed}

Before arriving at the separated solver, two other discretization
strategies were attempted and found inadequate.

\subsection{1D axial lattice}

The simplest possible molecular lattice places vertices along the
internuclear axis at $z \in \{0, R/4, R/2, 3R/4, R\}$ with nearest-
neighbor edges.  This fails catastrophically because:
\begin{enumerate}
  \item The 1D Coulomb potential $-Z/|z|$ is too singular without
    the $r^2$ Jacobian factor from 3D integration.
  \item The grid spacing $h = R/4$ artificially couples the kinetic
    energy scale to $R$.
  \item The 1D model has no physical correspondence to the 3D
    molecular problem.
\end{enumerate}
A cylinder-averaged variant (averaging the Coulomb potential over
a cylinder of radius $\rho_0$) produces qualitatively correct PES
shapes but requires $\rho_0$ as a free parameter, achieving only
$\sim 30\%$ energy error at best.

\subsection{2D tensor-product grid}

A direct 2D discretization on an $N_\xi \times N_\eta$ tensor product
grid in $(\xi, \eta)$ space solves the \emph{unseparated} problem
$H\Psi = E M \Psi$ where $M = \mathrm{diag}(\xi^2 - \eta^2)$ is
the volume weight.  While correct in principle, this approach suffers
from:
\begin{enumerate}
  \item Slow convergence: 28\% error at $50 \times 25 = 1250$
    unknowns, versus 1.01\% error at $N_\xi = 5000$ for the separated
    solver.
  \item Boundary condition difficulties: Dirichlet BCs at
    $\eta = \pm 1$ force $\Psi = 0$ at the nuclei, which is wrong
    for $\sigma_g$ states.  The correct natural zero-flux BCs are
    nontrivial to implement on a 2D grid.
  \item Computational cost: the 2D generalized eigenvalue problem
    is $O(N^3)$ in the total number of grid points, versus $O(N_\xi)$
    for the separated tridiagonal solver.
\end{enumerate}

The separated approach is $\sim 50\times$ more efficient because it
exploits the exact separability of the two-center Coulomb problem.

%% ========================================================================
%% BIBLIOGRAPHY
%% ========================================================================

\begin{thebibliography}{30}

\bibitem{loutey_paper0}
J.~Loutey,
``The Geometric Atom: Quantum State Space as a Packing Problem,''
GeoVac Paper~0 (2026).

\bibitem{loutey_paper1}
J.~Loutey,
``The Geometric Atom: Quantum Mechanics as a Packing Problem,''
GeoVac Paper~1 (2026).

\bibitem{loutey_paper7}
J.~Loutey,
``The Dimensionless Vacuum: Recovering the Schr\"odinger Equation
from Scale-Invariant Graph Topology,''
GeoVac Paper~7 (2026).

\bibitem{loutey_paper8}
J.~Loutey,
``Bond Sphere: Diatomic Molecules as $S^3$ and their Natural
Sturmian Basis,''
GeoVac Paper~8--9 (2026).

\bibitem{loutey_lcao_fci}
J.~Loutey,
``Topological LCAO Full CI for Heteronuclear Diatomics,''
GeoVac LCAO-FCI Paper (2026).

\bibitem{loutey_fci}
J.~Loutey,
``Full configuration interaction on a sparse graph Hamiltonian:
Multi-electron atoms via relational lattice indexing,''
GeoVac FCI Paper (2026).

\bibitem{loutey_paper10}
J.~Loutey,
``The Nuclear Lattice: Discrete Graph Structures for Molecular
Vibration and Rotation,''
GeoVac Paper~10 (2026).

\bibitem{Fock1935}
V.~Fock,
``Zur Theorie des Wasserstoffatoms,''
Z. Phys. \textbf{98}, 145--154 (1935).

\bibitem{Bates1953}
D.~R. Bates, K.~Ledsham, and A.~L. Stewart,
``Wave functions of the hydrogen molecular ion,''
Philos. Trans. R. Soc. London A \textbf{246}, 215--240 (1953).

\bibitem{Morse1953}
P.~M. Morse and H.~Feshbach,
\textit{Methods of Theoretical Physics}
(McGraw-Hill, New York, 1953).

\bibitem{Flammer1957}
C.~Flammer,
\textit{Spheroidal Wave Functions}
(Stanford University Press, Stanford, 1957).

\bibitem{James1933}
H.~M. James and A.~S. Coolidge,
``The ground state of the hydrogen molecule,''
J. Chem. Phys. \textbf{1}, 825--835 (1933).

\bibitem{Brent1973}
R.~P. Brent,
\textit{Algorithms for Minimization without Derivatives}
(Prentice-Hall, Englewood Cliffs, NJ, 1973).

\bibitem{Trefethen2000}
L.~N. Trefethen,
\textit{Spectral Methods in MATLAB}
(SIAM, Philadelphia, 2000).

\end{thebibliography}

\end{document}
